\documentclass[11pt,a4paper]{article}

\usepackage[utf8]{inputenc}

\usepackage[T1]{fontenc}

\usepackage[polish]{babel}

\usepackage{lmodern}

\usepackage{geometry}

\geometry{a4paper,margin=27mm}

\usepackage{microtype}

\usepackage[hidelinks]{hyperref}

% Portfolio: MSC-AR-P28; source: prospective ARC-XR research scope.

\title{Wyprzedzające wykrywanie ryzyka błędu manipulacji z kinematyki dłoni i VLM}

\author{Mateusz Pomianek}

\date{Wersja 0.3.0, 21 listopada 2026}

\begin{document}

\maketitle

\begin{abstract}

Przedmiotem opracowania jest Wyprzedzające wykrywanie ryzyka błędu manipulacji z kinematyki dłoni i VLM. Czy połączenie kinematyki dłoni i VLM wykrywa błędny chwyt przed kontaktem skuteczniej niż same trajektorie ruchu? Porównane zostaną modele oceniające ryzyko chwytu wyłącznie na podstawie kinematyki oraz modele dodatkowo korzystające z opisu sceny VLM. Głównym punktem końcowym jest czułość ostrzeżeń przed kontaktem przy stałym limicie fałszywych alarmów. Opis ma charakter wstępny i nie przedstawia niezweryfikowanych wyników.

\end{abstract}

\section{Wprowadzenie}

Rozwój adaptacyjnych systemów XR stwarza możliwość badania interakcji człowieka z informacją przestrzenną, ale samo zastosowanie AR, VR lub modeli sztucznej inteligencji nie przesądza o przewadze poznawczej czy ergonomicznej. W rozpatrywanym zagadnieniu, Wyprzedzające wykrywanie ryzyka błędu manipulacji z kinematyki dłoni i VLM, kluczowe jest rozdzielenie właściwości systemu od zmienności zachowania użytkownika. Porównanie należy zaprojektować tak, aby mierzyć jednoznaczne zjawisko, a nie tylko ogólną ocenę technologii.

\section{Cel pracy i pytanie badawcze}

Czy połączenie kinematyki dłoni i VLM wykrywa błędny chwyt przed kontaktem skuteczniej niż same trajektorie ruchu? Projekt wymaga krytycznego przeglądu aktualnego stanu badań oraz zdefiniowania nowości w stosunku do innych tematów portfolio XR.

\section{Zakres i odrębny wkład naukowy}
Przedkontaktowa predykcja niepoprawnego chwytu na podstawie kinematyki dłoni i opisu relacji obiektów z VLM; nie jest to identyfikacja obiektu 3D (P04) ani wykrywanie brakującego kroku po wykonaniu (P05). Rzeczywista prognoza ryzyka manipulacji przed zdarzeniem, a nie wyjaśnienie błędu po wykonaniu kroku.
\section{Hipotezy badawcze}
\paragraph{H1.} Fuzja VLM i kinematyki zwiększa czułość wczesnej predykcji błędnego chwytu przy tym samym budżecie fałszywych alarmów.
\paragraph{H2.} Korzyść dotyczy chwytów w scenach z niejednoznacznym obiektem i relacją, a nie wyłącznie zadań o jednym celu.
\paragraph{H3.} Zysk utrzymuje się przy uwzględnieniu pełnego opóźnienia VLM, tak aby alarm rzeczywiście poprzedzał kontakt.
\section{Projekt badania i plan eksperymentów}
\subsection{A. Operacjonalizacja i kwalifikacja danych}
Przedkontaktowa predykcja niepoprawnego chwytu na podstawie kinematyki dłoni i opisu relacji obiektów z VLM; nie jest to identyfikacja obiektu 3D (P04) ani wykrywanie brakującego kroku po wykonaniu (P05). Jednostka wnioskowania: niezależna próba chwytu z oznaczonym pierwszym kontaktem i poprawnością. Przed badaniem należy określić źródła, wersje urządzeń, czas akwizycji, etykiety i niezależną referencję. Nie liczyć zależnych okien i klatek jako niezależnych prób.
\subsection{B. Eksperyment porównawczy}
Niezależnie zaetykietować obiekt, trajektorię dłoni, początek kontaktu i błąd chwytu; podać zamrożony model kinematyczny i wariant VLM. Wynik pierwszorzędowy: różnica czułości ostrzeżeń przed kontaktem przy stałej liczbie fałszywych alarmów na minutę próby między VLM+kinematyką a samą kinematyką. Zamrozić losowanie warunków, progi, mianownik punktu końcowego i kryteria wyłączeń bez dostępu do oceny końcowej.
\subsection{C. Kontrole i analiza odporności}
Bez dostępu do obrazu po momencie predykcji, z jawnym ograniczeniem opóźnienia obliczeniowego i niepewności głębi. Ocenić fałszywe alarmy, utratę śledzenia dłoni i błąd segmentacji, oraz podział po osobie i obiekcie. Utrata sygnału, abstencja, przerwanie zadania i dane brakujące są odrębnymi kategoriami wyniku.
\subsection{D. Walidacja na niezależnym materiale}
Walidacja musi zachować podział po uczestniku, sesji, scenie lub zadaniu. Krzywe czułości i lead time przy stałych false alerts/min z bootstrapem próba/chwyt, analiza opóźnienia end-to-end. Raportować wyniki negatywne i rozbieżności z założeniami.
\section{Estymand, analiza statystyczna i decyzje}
różnica czułości ostrzeżeń przed kontaktem przy stałej liczbie fałszywych alarmów na minutę próby między VLM+kinematyką a samą kinematyką. Jednostka: niezależna próba chwytu z oznaczonym pierwszym kontaktem i poprawnością. Krzywe czułości i lead time przy stałych false alerts/min z bootstrapem próba/chwyt, analiza opóźnienia end-to-end. H1 jest planowanym kontrastem głównym; H2 i H3 pozostają pomocnicze do zatwierdzenia hierarchii testów oraz kontroli wielokrotności. Wartości marginesów, mocy, minimalnego efektu istotnego praktycznie i liczebności mają wynikać z pilotażu, a nie z post hoc wyboru wyniku.
\section{Stan badań i granica nowości}
\begin{itemize}
\item 3D grounding with VLM is prior art; the new study must demonstrate calibrated AR anchor gains, not merely combining 2D and 3D. \cite{xrprior1}.
\item Prior hand-gesture AR guidance has tested error feedback; establish a distinct motor-phase or procedural endpoint. \cite{xrprior2}.
\end{itemize}
Rzeczywista prognoza ryzyka manipulacji przed zdarzeniem, a nie wyjaśnienie błędu po wykonaniu kroku. Odrębność na tle opublikowanych metod i pozostałych artykułów wymaga osobnego audytu.
\section{Warunki bezpieczeństwa, sprzętu i zgodności}
To prospektywna koncepcja, nie zgoda na eksperyment, prerejestracja lub wynik empiryczny. Praca z uczestnikami, interwencje stresowe i pomiar biometryczny wymagają przeglądu etycznego, świadomej zgody i ochrony danych. Quest 3 w wersji standardowej nie posiada eye trackingu, a Tobii Eye Tracker 5 jest urządzeniem ekranowym; nie zastępować jednego stanowiska drugim.
\section{Materiał i metody}

Porównane zostaną modele oceniające ryzyko chwytu wyłącznie na podstawie kinematyki oraz modele dodatkowo korzystające z opisu sceny VLM. Wymagane jest zdefiniowanie jednostki analizy, zasad doboru przypadków, wersji scenariusza i źródeł pomiaru. Rozkład warunków oraz reguły wykluczenia prób należy określić przed pomiarem. W badaniach z udziałem ludzi niezbędna jest uprzednia właściwa ocena etyczna i świadoma zgoda. W testach wyłącznie symulacyjnych należy oddzielić prawdę symulacyjną od obserwacji empirycznej.

\subsection{Wynik główny i kryteria jakości}

Czułość ostrzeżeń przed kontaktem przy stałym limicie fałszywych alarmów. Pomiar powinien zachowywać jawne jednostki, znaczniki czasu i reguły oceny poprawności. Przy modelach predykcyjnych stosować rozłączne partycje oceny; przy eksperymentach interwencyjnych oddzielić skuteczność predykcji od przyczynowego wpływu interfejsu.

\section{Plan analizy}

Analiza czasu wyprzedzenia, czułości warunkowej i uogólniania na nowe obiekty i uczestników. W wynikach należy raportować wielkość efektu lub błąd estymacji z przedziałem niepewności, liczbą obserwacji i odsetkiem odrzuconych prób. Analizy dodatkowe wymagają jasnego oznaczenia jako eksploracyjne. Nie wybierać progów decyzyjnych ani modelu na podstawie obserwacji ze zbioru testowego.

\section{Wyniki}

W obecnej wersji nie raportuje się wyników pomiarów, tabel statystycznych ani parametrów wyznaczonych z nieudokumentowanych danych. To kierunek badawczy, a nie zatwierdzony protokół badania ani zgłoszony artykuł.

\section{Dyskusja}

Interpretacja przyszłych wyników musi uwzględnić alternatywne mechanizmy i zależność wniosków od jakości danych, warunków zadania oraz ograniczeń sprzętu. Wartość potencjalnego rozwiązania zależy nie tylko od trafności modelu, lecz także od jego opóźnienia, stabilności i wpływu na zachowanie człowieka. Wyniki testów syntetycznych nie mogą być automatycznie uogólniane na populację uczestników.

\section{Ograniczenia badania}

Niejednoznaczna segmentacja dłoni oraz opóźnienie VLM mogą przekreślić praktyczną przewagę modelu. Predykcja ryzyka przed wykonaniem czynności, nie przestrzenne przypisanie VLM MSC-AR-P04 lub detekcja pominiętych kroków MSC-AR-P05. Żadna część szkicu nie dowodzi przeprowadzenia zbierania danych, uzyskania zgód etycznych ani gotowości manuskryptu do zgłoszenia.

\section{Wnioski}

Przedstawiono problem, kontrast i sposób jego planowanej oceny. Rozstrzygnięcie postawionego pytania wymaga przygotowania protokołu, zweryfikowanych źródeł oraz poprawnie udokumentowanych wyników. Bez tych materiałów nie można sformułować konkluzji empirycznych.

\begin{thebibliography}{9}
\bibitem{xrprior1} SeeGround: See and Ground for Zero-Shot Open-Vocabulary 3D Visual Grounding (2025), CVPR 2025. \url{https://ieeexplore.ieee.org/abstract/document/11094047}.
\bibitem{xrprior2} Exploring AR hand augmentations as error feedback mechanisms for enhancing gesture-based tutorials (2025), Frontiers in Virtual Reality. \url{https://doi.org/10.3389/frvir.2025.1574965}.
\end{thebibliography}
\end{document}

